\section{What we take away from this}
\label{sec:takeaways}

\paragraph{For the research story.} The interesting result is not ``LLM beats
baselines'' --- it is the shape of the robustness: no cliffs under thinning, real
costs only under outright loss or bad representation, and a deep data layer whose
value shows up as investigation speed and explanation quality rather than headline
accuracy. Those are claims none of the existing automated methods could even
express, because they cannot use the kernel layer at all and their accuracy is flat
for the uninteresting reason that it is low.

\paragraph{For system builders.} Three transferable design lessons: represent deep
data before showing it to a model (interpretation beats access); make tools answer
definitively when a data source is absent, or the agent wastes a third of its steps;
and a deterministic evidence briefing bought more accuracy-per-dollar than any
prompt engineering we tried.

\paragraph{For observability practice.} Teams over-collect for this use case.
Five-percent trace sampling, one-minute metrics and error-only logs kept over 90\%
of localization ability in our setting; the kernel layer is cheap to record
($-0.5\%$ throughput) and repays itself in faster, better-explained diagnoses.

\paragraph{For evaluation practice.} Rigour changed our conclusions twice: hiding
answer-revealing identifiers collapsed early (inflated) results and forced us to
build genuinely better tools; and measuring run-to-run noise stopped us from
over-reading single-run differences. Both disciplines --- plus full recording of
every diagnosis --- are what make the numbers in this report trustworthy enough to
discuss.
